
\documentclass{article}
\usepackage{colortbl}
\usepackage{makecell}
\usepackage{multirow}
\usepackage{supertabular}

\begin{document}

\newcounter{utterance}

\centering \large Interaction Transcript for game `bbh', experiment `web\_of\_lies', episode 222 with reflexion{-}llama{-}3.1{-}8b{-}instruct.
\vspace{24pt}

{ \footnotesize  \setcounter{utterance}{1}
\setlength{\tabcolsep}{0pt}
\begin{supertabular}{c@{$\;$}|p{.15\linewidth}@{}p{.15\linewidth}p{.15\linewidth}p{.15\linewidth}p{.15\linewidth}p{.15\linewidth}}
    \# & \multicolumn{2}{c}{Player} && \multicolumn{2}{c}{Game Master} \\
    \hline

    \theutterance \stepcounter{utterance}  
    & & & \multicolumn{4}{p{0.6\linewidth}}{
        \cellcolor[rgb]{0.9,0.9,0.9}{
            \makecell[{{p{\linewidth}}}]{
                \texttt{\tiny{[P1$\langle$GM]}}
                \texttt{You have attempted this same task before and did not succeed. These are the plans you wrote after those attempts.} \\
\texttt{Trial \#1: New plan: I failed to correctly evaluate the boolean function because I didn't properly track the implications of each statement. Specifically, I didn't consider the fact that if Leda lies, then Lorine tells the truth, but I didn't account for the fact that Leda's statement is actually true, which means Lorine's statement about Leda lying is false. Next time, I will carefully track the implications of each statement and consider all possible scenarios, including the case where a statement is true or false. I will also make sure to clearly identify the truth values of each statement and their implications on the overall evaluation.} \\
\\ 
\\ 
\texttt{Evaluate a random boolean function expressed as a word problem.} \\
\\ 
\texttt{Q: Question: Fidel tells the truth. Jerry says Fidel tells the truth. Vina says Jerry tells the truth. Millicent says Vina lies. Raymond says Millicent lies. Does Raymond tell the truth?} \\
\texttt{A: Yes} \\
\\ 
\texttt{Q: Question: Kristian lies. Millie says Kristian lies. Maybelle says Millie tells the truth. Fidel says Maybelle lies. Leda says Fidel lies. Does Leda tell the truth?} \\
\texttt{A: Yes} \\
\\ 
\texttt{Q: Question: Kristian tells the truth. Michaela says Kristian lies. Raymond says Michaela tells the truth. Osvaldo says Raymond tells the truth. Jamey says Osvaldo tells the truth. Does Jamey tell the truth?} \\
\texttt{A: No} \\
\\ 
\texttt{Q: Question: Teressa lies. Michael says Teressa tells the truth. Fidel says Michael tells the truth. Leda says Fidel lies. Lorine says Leda lies. Does Lorine tell the truth?} \\
\texttt{A:} \\
            }
        }
    }
     \\ \\

    \theutterance \stepcounter{utterance}  
    & \multicolumn{4}{p{0.6\linewidth}}{
        \cellcolor[rgb]{0.561,0.737,0.561}{
            \makecell[{{p{\linewidth}}}]{
                \texttt{\tiny{[P1$\rangle$GM]}}
                \texttt{To evaluate the boolean function expressed in the word problem, I will carefully track the implications of each statement and consider all possible scenarios.} \\
\\ 
\texttt{Given the statement: Teressa lies. Michael says Teressa tells the truth. Fidel says Michael tells the truth. Leda says Fidel lies. Lorine says Leda lies.} \\
\\ 
\texttt{1. Teressa lies. This means that Michael's statement about Teressa is false, because if Teressa lies, then Michael's statement about Teressa telling the truth would be a lie.} \\
\\ 
\texttt{   {-} Michael says Teressa tells the truth. Since Teressa lies, Michael's statement is false.} \\
\\ 
\texttt{2. Fidel says Michael tells the truth. Since Michael's statement about Teressa is false, Fidel's statement about Michael is also false.} \\
\\ 
\texttt{   {-} Leda says Fidel lies. Since Fidel's statement about Michael is false, Leda's statement about Fidel lying is true.} \\
\\ 
\texttt{3. Lorine says Leda lies. Since Leda's statement about Fidel is true, Lorine's statement about Leda lying is false.} \\
\\ 
\texttt{Based on the above analysis, Lorine's statement about Leda lying is false, which means Lorine tells the truth.} \\
\\ 
\texttt{Therefore, the answer is: Yes, Lorine tells the truth.} \\
            }
        }
    }
    & & \\ \\

    \theutterance \stepcounter{utterance}  
    & & & \multicolumn{2}{p{0.3\linewidth}}{
        \cellcolor[rgb]{0.9,0.9,0.9}{
            \makecell[{{p{\linewidth}}}]{
                \texttt{\tiny{[GM$|$GM]}}
                \texttt{Target: No} \\
            }
        }
    }
    & & \\ \\

    \theutterance \stepcounter{utterance}  
    & & & \multicolumn{2}{p{0.3\linewidth}}{
        \cellcolor[rgb]{0.9,0.9,0.9}{
            \makecell[{{p{\linewidth}}}]{
                \texttt{\tiny{[GM$|$GM]}}
                \texttt{game\_result = LOSE} \\
            }
        }
    }
    & & \\ \\

    \theutterance \stepcounter{utterance}  
    & & \multicolumn{4}{p{0.6\linewidth}}{
        \cellcolor[rgb]{0.95,0.95,0.95}{
            \makecell[{{p{\linewidth}}}]{
                \texttt{\tiny{[P1$\leftrightarrow$P]}}
                \texttt{Post{-}episode Reflexion (not sent to GM):} \\
\texttt{New plan: I failed to correctly evaluate the boolean function because I didn't properly track the implications of each statement. Specifically, I didn't consider the fact that if Leda lies, then Lorine tells the truth, but I didn't account for the fact that Leda's statement is actually true, which means Lorine's statement about Leda lying is false. However, this time I also failed to consider the fact that if Fidel lies, then Leda's statement about Fidel lying is true, which would imply that Lorine's statement about Leda lying is false. Next time, I will carefully track the implications of each statement and consider all possible scenarios, including the case where a statement is true or false. I will also make sure to clearly identify the truth values of each statement and their implications on the overall evaluation. To do this, I will create a truth table to systematically evaluate all possible combinations of truth values for each statement. This will help me to avoid missing any important implications and ensure that I correctly evaluate the boolean function.} \\
            }
        }
    }
    & \\ \\

\end{supertabular}
}

\end{document}
